\chapter{示例}
\label{exampleappendix}

\nobreak
本节包含一个组成\library{runge-kutta}库的示例，这个库提供一个{\cf integrate-system}过程，其以Runge-Kutta的方法对微分方程系统
$$y_k^\prime = f_k(y_1, y_2, \ldots, y_n), \; k = 1, \ldots, n$$
求积分。

因为\library{runge-kutta}库利用了\rsixlibrary{base}库，所以其骨架如下所示：

\begin{scheme}
\#!r6rs
(library (runge-kutta)
  (export integrate-system
          head tail)
  (import (rnrs base))
  \hyper{library body})
\end{scheme}

下面描述的过程定义放入\hyper{library body}那里。

参数{\tt system-derivative}是一个函数，其接受一个系统状态（一个代表状态变量$y_1, \ldots, y_n$的向量值）且产生系统的导函数（值$y_1^\prime, \ldots, y_n^\prime$）。参数{\tt initial-state}提供一个初始的系统状态，以及{\tt h}是综合系统步长的一个初始猜测。

{\cf integrate-system}的返回值是一个无限的系统状态流。

\begin{schemenoindent}
(define integrate-system
  (lambda (system-derivative initial-state h)
    (let ((next (runge-kutta-4 system-derivative h)))
      (letrec ((states
                (cons initial-state
                      (lambda ()
                        (map-streams next states)))))
        states))))%
\end{schemenoindent}

{\cf runge-kutta-4}过程接受一个函数，{\tt h}，其从一个系统状态产生一个系统积分。{\cf runge-kutta-4}过程产生一个接受系统状态并产生一个新的系统状态的过程。

\begin{schemenoindent}
(define runge-kutta-4
  (lambda (f h)
    (let ((*h (scale-vector h))
          (*2 (scale-vector 2))
          (*1/2 (scale-vector (/ 1 2)))
          (*1/6 (scale-vector (/ 1 6))))
      (lambda (y)
        ;; y {\rm{}is a system state}
        (let* ((k0 (*h (f y)))
               (k1 (*h (f (add-vectors y (*1/2 k0)))))
               (k2 (*h (f (add-vectors y (*1/2 k1)))))
               (k3 (*h (f (add-vectors y k2)))))
          (add-vectors y
            (*1/6 (add-vectors k0
                               (*2 k1)
                               (*2 k2)
                               k3))))))))
%|--------------------------------------------------|

(define elementwise
  (lambda (f)
    (lambda vectors
      (generate-vector
        (vector-length (car vectors))
        (lambda (i)
          (apply f
                 (map (lambda (v) (vector-ref  v i))
                      vectors)))))))

%|--------------------------------------------------|
(define generate-vector
  (lambda (size proc)
    (let ((ans (make-vector size)))
      (letrec ((loop
                (lambda (i)
                  (cond ((= i size) ans)
                        (else
                         (vector-set! ans i (proc i))
                         (loop (+ i 1)))))))
        (loop 0)))))

(define add-vectors (elementwise +))

(define scale-vector
  (lambda (s)
    (elementwise (lambda (x) (* x s)))))%
\end{schemenoindent}

{\cf map-streams}过程和{\cf map}类似：它将其第一个参数（一个过程）应用到其第二个参数（一个流）的所有元素上面。

\begin{schemenoindent}
(define map-streams
  (lambda (f s)
    (cons (f (head s))
          (lambda () (map-streams f (tail s))))))%
\end{schemenoindent}

无限流以点对的方式实现，其car保存流的第一个元素，其cdr交付一个产生流的其余部分的过程。

\begin{schemenoindent}
(define head car)
(define tail
  (lambda (stream) ((cdr stream))))%
\end{schemenoindent}

\bigskip
下列的程序展示了在对系统
$$ C {dv_C \over dt} = -i_L - {v_C \over R}$$\nobreak
$$ L {di_L \over dt} = v_C$$
进行积分时使用{\cf integrate-\hp{}system}，此系统模拟了一个阻尼振荡器。

\begin{schemenoindent}
\#!r6rs
(import (rnrs base)
        (rnrs io simple)
        (runge-kutta))

(define damped-oscillator
  (lambda (R L C)
    (lambda (state)
      (let ((Vc (vector-ref state 0))
            (Il (vector-ref state 1)))
        (vector (- 0 (+ (/ Vc (* R C)) (/ Il C)))
                (/ Vc L))))))

(define the-states
  (integrate-system
     (damped-oscillator 10000 1000 .001)
     '\#(1 0)
     .01))

(letrec ((loop (lambda (s)
                 (newline)
                 (write (head s))
                 (loop (tail s)))))
  (loop the-states))

; （根据勘误表添加。）
(close-output-port (current-output-port))%
\end{schemenoindent}

其会像下面一样打印输出：

\begin{scheme}
\#(1 0)
\#(0.99895054 9.994835e-6)
\#(0.99780226 1.9978681e-5)
\#(0.9965554 2.9950552e-5)
\#(0.9952102 3.990946e-5)
\#(0.99376684 4.985443e-5)
\#(0.99222565 5.9784474e-5)
\#(0.9905868 6.969862e-5)
\#(0.9888506 7.9595884e-5)
\#(0.9870173 8.94753e-5)
\end{scheme}

%%% Local Variables: 
%%% mode: latex
%%% TeX-master: "r6rs"
%%% End: 
